% 作者 shingja · github.com/shingja/shingja-math-modeling-workbench · 双许可见 LICENSE · 保留版权标识
% ============================================================
% 02 — 模型假设
% 说明: 列出建模所做的合理简化假设, 每条编号。
%       假设应具体、可验证, 并简要说明理由。
%       数值必须从 results/results.json 提取, 禁止手写。
% ============================================================

\section{模型假设}

为简化问题并建立可求解的数学模型，本文做出以下合理假设：

\begin{enumerate}
    \item \textbf{[假设名称]:} [假设内容及合理性说明]
    \item \textbf{[假设名称]:} [假设内容及合理性说明]
    \item \textbf{[假设名称]:} [假设内容及合理性说明]
    \item \textbf{[假设名称]:} [假设内容及合理性说明]
    \item \textbf{[假设名称]:} [假设内容及合理性说明]
\end{enumerate}

以上假设在简化模型的同时保留了问题的核心特征，后续章节将在此基础上展开建模。
